% vim: tw=78 ai
\documentclass{article}

%\usepackage{a4wide}
\usepackage{graphicx}
\usepackage[utf8x]{inputenc}
\usepackage[czech]{babel}
\usepackage[IL2]{fontenc}

\pagestyle{empty}

\begin{document}
\begin{flushright}
\today
\end{flushright}
\begin{center}
\Large\textbf{Posudek vedoucího na bakalářskou práci \\
\normalsize%
\textit{Erich Duda:} Nástroje pro vyhledávání definic a použití v jazyce C}
\end{center}

Dle zadání měla bakalářská práce zmapovat nástroje pro vyhledávání definic a
použití proměnných v programech psaných jazykem C. Text práce je na 41
stranách strukturovaný do 6 kapitol. V úvodu autor diskutuje motivaci pro
práci, na kterou dále navazuje analýzou současných nástrojů a jejich
nedostatků, shrnujíc vše na konci kapitoly. Dále se autor věnuje teoretickému
základu nutnému pro vývoj nástroje, který odstraní výše uvedené nedostatky.
Popis samotného vývoje se odehrává v další kapitole s vysvětlením zádrhelů a
navrženými řešeními. Následuje kapitola, která dokumentuje použití nástroje.
Práci ukončuje závěr s popisem práce do budoucna.

\textbf{Rozsah} uvedených kapitol je z mého pohledu dostatečný.
\textbf{Obsahově} jsem neshledal nedostatky, práce obsahuje vše, co dle zadání
měla. Po \textbf{jazykové} stránce taktéž neshledávám práci nijak zanedbanou,
neobsahuje překlepy ani jiné jazykové deviace. Možná trochu ruší příliš velké
odsazení kolem obrázků (např. str. 34) a úvody podkapitol mohly být zvládnuté
lépe.

\textbf{Praktická část} práce je v přiloženém \texttt{zip} souboru. Kód je
dobře čitelný a funkční, nicméně není v některých místech zcela bezpečný.
Např. metoda \texttt{isCFile} přistupuje mimo pole, bude-li vstupem soubor s
názvem ,,c''.

Na závěr konstatuji, že práce \textbf{splňuje zadání a požadavky na
bakalářskou práci}. Navrhuji práci \textbf{uznat} jako bakalářskou, a vzhledem k
výše uvedeným faktům doporučuji hodnocení stupněm B.

\vspace{3em}
\begin{flushright}
Jiří Slabý
\end{flushright}

\end{document}
